\documentclass[11pt]{article}

\usepackage[a4paper,margin=1in]{geometry}
\usepackage{amsmath,amssymb,amsthm,mathtools}
\usepackage{booktabs}
\usepackage{microtype}
\usepackage[hidelinks]{hyperref}
\usepackage{aliascnt}
\usepackage[nameinlink,noabbrev]{cleveref}

\allowdisplaybreaks

\newtheorem{theorem}{Theorem}[section]
\newaliascnt{proposition}{theorem}
\newtheorem{proposition}[proposition]{Proposition}
\aliascntresetthe{proposition}
\newaliascnt{lemma}{theorem}
\newtheorem{lemma}[lemma]{Lemma}
\aliascntresetthe{lemma}
\newaliascnt{corollary}{theorem}
\newtheorem{corollary}[corollary]{Corollary}
\aliascntresetthe{corollary}
\theoremstyle{definition}
\newaliascnt{definition}{theorem}
\newtheorem{definition}[definition]{Definition}
\aliascntresetthe{definition}
\theoremstyle{remark}
\newaliascnt{remark}{theorem}
\newtheorem{remark}[remark]{Remark}
\aliascntresetthe{remark}

\newcommand{\R}{\mathbb{R}}
\newcommand{\dist}{\operatorname{dist}}
\newcommand{\ind}{\mathbf{1}}

\title{Distinct Distances in Planar Point Sets\\with No Four Concyclic Points}
\author{Jamal Agbanwa}
\date{July 2026}

\hypersetup{
  pdftitle={Distinct Distances in Planar Point Sets with No Four Concyclic Points},
  pdfauthor={Jamal Agbanwa},
  pdfsubject={Bounds for distinct distances under a no-four-concyclicity condition},
  pdfkeywords={distinct distances, concyclic points, diameter graph, extremal geometry}
}

\begin{document}

\maketitle

\begin{abstract}
For a set $P$ of $n$ distinct points in the Euclidean plane, let $d_P(p)$ denote the number of distinct distances from $p\in P$ to the other points of $P$.  We study the least possible value of $\max_{p\in P}d_P(p)$ under the restriction that no four points of $P$ lie on one circle.  A direct multiplicity argument gives only $\lceil(n-1)/3\rceil$.  We show, using an elementary bound for the number of diameter pairs in a planar set, that the exceptional congruence class can also be improved, yielding
\[
   f(n)\ge \left\lceil\frac n3\right\rceil.
\]
We then give an explicit construction on two perpendicular axes proving
\[
   f(n)\le \left\lfloor\frac{3n}{4}\right\rfloor
          -\ind_{\{4\mid n\}}.
\]
Consequently $\limsup f(n)/n\le 3/4$, so an asymptotic lower bound of the form $f(n)>(1-o(1))n$ is impossible.  The note is completely self-contained.  We also formulate exactly what additional ``distance-fibre deficiency'' would be required to improve the lower bound to $(1/3+c)n$.
\end{abstract}

\medskip
\noindent\textbf{Keywords.} Distinct distances; concyclic points; diameter graph; extremal planar geometry.

\section{Definition and main result}

Throughout, a point configuration means a finite \emph{set} of pairwise distinct points.  Thus the condition that no four points are concyclic always refers to four distinct points and to an ordinary Euclidean circle.

\begin{definition}
Let $P\subset\R^2$ be finite and let $p\in P$.  Define
\[
 d_P(p):=\left|\bigl\{|p-q|:q\in P\setminus\{p\}\bigr\}\right|.
\]
For $n\ge 1$, let
\[
 f(n):=
 \min_{\substack{P\subset\R^2,\ |P|=n\\
                  \text{no four points of $P$ are concyclic}}}
       \ \max_{p\in P}d_P(p).
\]
\end{definition}

The minimum is well defined: the class of admissible configurations is nonempty, and the quantity being minimized is an integer in $\{0,1,\ldots,n-1\}$.

\begin{theorem}[Main bounds]\label{thm:main}
For every integer $n\ge 2$,
\[
 \boxed{
 \left\lceil\frac n3\right\rceil
 \le f(n)
 \le
 \left\lfloor\frac{3n}{4}\right\rfloor-
 \ind_{\{4\mid n\}}.}
\]
Equivalently, if $m\ge 0$, then the upper bound is
\[
 f(n)\le
 \begin{cases}
  3m-1,&n=4m\quad(m\ge1),\\
  3m,&n=4m+1,\\
  3m+1,&n=4m+2,\\
  3m+2,&n=4m+3.
 \end{cases}
\]
\end{theorem}

Here $\ind_{\{4\mid n\}}$ denotes the indicator of the divisibility condition $4\mid n$: it is $1$ when $4$ divides $n$ and $0$ otherwise.

\begin{corollary}\label{cor:asymptotic}
One has
\[
 \frac13
 \le \liminf_{n\to\infty}\frac{f(n)}n
 \le \limsup_{n\to\infty}\frac{f(n)}n
 \le \frac34.
\]
In particular, there is no sequence $\varepsilon_n\to0$ such that
\[
 f(n)>(1-\varepsilon_n)n
\]
for all sufficiently large $n$.  Thus the proposed estimate
\[
 f(n)>(1-o(1))n
\]
is false.
\end{corollary}

\begin{proof}
The limit inequalities follow immediately from \cref{thm:main}.  The last assertion follows because the upper bound in \cref{thm:main} gives $f(n)/n\le 3/4$ for every $n\ge2$, whereas $(1-\varepsilon_n)>3/4$ for all sufficiently large $n$ whenever $\varepsilon_n\to0$.
\end{proof}

The lower and upper bounds are proved separately in \cref{sec:lower,sec:upper}.

\section{The elementary distance-fibre bound}\label{sec:fibres}

Fix an admissible configuration $P$ of cardinality $n$ and a point $p\in P$.  For each distance $r>0$ realized from $p$, the corresponding distance fibre is
\[
 F_p(r):=\{q\in P\setminus\{p\}:|p-q|=r\}.
\]
It lies on the circle of centre $p$ and radius $r$.  Hence the hypothesis immediately controls its cardinality.

\begin{proposition}[Capacity bound]\label{prop:capacity}
For every $p\in P$ and every realized distance $r>0$,
\[
 |F_p(r)|\le3.
\]
Consequently
\[
 d_P(p)\ge \left\lceil\frac{n-1}{3}\right\rceil.
\]
\end{proposition}

\begin{proof}
If $|F_p(r)|\ge4$, then four distinct points of $P$ lie on the circle with centre $p$ and radius $r$, contrary to the hypothesis.  The fibres $F_p(r)$ over the $d_P(p)$ realized distances partition the $n-1$ points of $P\setminus\{p\}$.  Therefore
\[
 n-1=\sum_r |F_p(r)|\le3d_P(p),
\]
which gives the stated integer lower bound.
\end{proof}

For $n\not\equiv1\pmod3$, this already equals $\lceil n/3\rceil$.  If $n=3m+1$, however, it gives only $m$ rather than $m+1$.  The missing unit is forced by the geometry of diameter pairs.

\section{A self-contained diameter-graph bound}\label{sec:diameter}

Let $S\subset\R^2$ be a finite set with at least two points, and write
\[
 D(S):=\max\{|x-y|:x,y\in S\}.
\]
The \emph{diameter graph} $G_D(S)$ has vertex set $S$, with an edge $xy$ precisely when $|x-y|=D(S)$.

We shall prove the planar bound
\[
 |E(G_D(S))|\le |S|.
\]
The proof is included in full because this is the only genuinely geometric ingredient needed for the lower bound.

\begin{lemma}[Angular concentration of diameter neighbours]\label{lem:angular}
Let $x\in S$, and let $y,z\in S\setminus\{x\}$ satisfy
\[
 |x-y|=|x-z|=D(S).
\]
Then
\[
 \angle yxz\le\frac{\pi}{3}.
\]
Consequently, all rays from $x$ to its neighbours in $G_D(S)$ are contained in some closed angular arc of width at most $\pi/3$.
\end{lemma}

\begin{proof}
Put $D=D(S)$ and $\theta=\angle yxz\in[0,\pi]$.  By the law of cosines,
\[
 |y-z|^2=D^2+D^2-2D^2\cos\theta
          =2D^2(1-\cos\theta).
\]
Since $D$ is the diameter of $S$, one has $|y-z|\le D$.  Hence
\[
 2(1-\cos\theta)\le1,
\]
so $\cos\theta\ge1/2$ and therefore $\theta\le\pi/3$.

It remains to justify the final assertion, because pairwise angular separation is a circular rather than a linear condition.  Choose one diameter-neighbour direction as angular coordinate $0$.  Every other neighbour direction has a representative in $[-\pi/3,\pi/3]$.  Let $a$ and $b$ be the minimum and maximum of these representatives.  If $a$ and $b$ have the same sign, then plainly $b-a\le\pi/3$.  If $a<0<b$, then
\[
 b-a\le\frac{2\pi}{3}<\pi.
\]
Thus $b-a$ is the smaller angular separation between the two extreme directions; by the first part of the proof it is at most $\pi/3$.  Hence all neighbour directions lie in the closed arc represented by $[a,b]$, whose width is at most $\pi/3$.
\end{proof}

Fix once and for all an orientation of the plane, so that ``clockwise'' has an unambiguous meaning.  For each non-isolated vertex $x$ of $G_D(S)$, choose a closed arc of width at most $\pi/3$ containing all its diameter-neighbour directions, as supplied by \cref{lem:angular}.  Let $c(x)$ be the unique diameter-neighbour whose ray is clockwise-most in that arc.  Uniqueness holds because two distinct points at the same positive distance from $x$ cannot lie on the same ray from $x$.

\begin{lemma}[Every diameter edge is selected at an endpoint]\label{lem:selected}
If $xy$ is an edge of $G_D(S)$, then
\[
 c(x)=y\qquad\text{or}\qquad c(y)=x.
\]
\end{lemma}

\begin{proof}
Suppose for contradiction that neither equality holds.  Since $y$ is not clockwise-most among the diameter-neighbours of $x$, there is a diameter-neighbour $y'$ of $x$ obtained from the ray $xy$ by a strictly positive clockwise turn through some angle
\[
 0<\alpha\le\frac{\pi}{3}.
\]
Similarly, there is a diameter-neighbour $x'$ of $y$ obtained from the ray $yx$ by a strictly positive clockwise turn through an angle
\[
 0<\beta\le\frac{\pi}{3}.
\]

Apply a direct similarity of the plane, consisting of a translation, a rotation, and a scaling, so that
\[
 x=(0,0),\qquad y=(1,0),
\]
and the common diameter becomes $1$.  A clockwise turn from the positive horizontal ray places $y'$ below the $x$-axis, whereas a clockwise turn from the negative horizontal ray places $x'$ above it.  Thus
\[
 y'=(\cos\alpha,-\sin\alpha),
 \qquad
 x'=(1-\cos\beta,\sin\beta).
\]
It follows that
\begin{align*}
 |y'-x'|^2-1
 &= (\cos\alpha+\cos\beta-1)^2
    +(\sin\alpha+\sin\beta)^2-1\\
 &=4\cos\frac{\alpha-\beta}{2}
   \left(\cos\frac{\alpha-\beta}{2}
      -\cos\frac{\alpha+\beta}{2}
   \right).
\end{align*}
Now
\[
 \left|\frac{\alpha-\beta}{2}\right|
 <\frac{\alpha+\beta}{2}\le\frac{\pi}{3}.
\]
The cosine function is even and strictly decreasing on $[0,\pi]$.  Therefore
\[
 \cos\frac{\alpha-\beta}{2}>0
 \quad\text{and}\quad
 \cos\frac{\alpha-\beta}{2}
 >\cos\frac{\alpha+\beta}{2}.
\]
Hence $|y'-x'|^2-1>0$, so $|y'-x'|>1$.  This contradicts the fact that the diameter is $1$.
\end{proof}

\begin{theorem}[Planar diameter bound]\label{thm:diameter}
Every finite set $S\subset\R^2$ with at least two points determines at most $|S|$ diameter pairs:
\[
 |E(G_D(S))|\le |S|.
\]
\end{theorem}

\begin{proof}
By \cref{lem:selected}, every diameter edge $xy$ has at least one endpoint at which it is the selected edge: either $c(x)=y$ or $c(y)=x$.  Assign $xy$ to one such endpoint, choosing arbitrarily if both endpoints select it.  A vertex $x$ can receive at most one edge, namely the single edge $xc(x)$.  Thus the assignment is an injection from the edge set into the vertex set, and
\[
 |E(G_D(S))|\le|S|.
\]
\end{proof}

\section{The lower bound}\label{sec:lower}

\begin{theorem}[Lower bound]\label{thm:lower}
Let $P\subset\R^2$ be an admissible configuration of $n\ge2$ points.  Then some point $p\in P$ determines at least $\lceil n/3\rceil$ distinct distances.  Equivalently,
\[
 f(n)\ge\left\lceil\frac n3\right\rceil.
\]
\end{theorem}

\begin{proof}
By \cref{prop:capacity}, every point $p\in P$ satisfies
\[
 d_P(p)\ge\left\lceil\frac{n-1}{3}\right\rceil.
\]
If $n=3m$ or $n=3m+2$, then
\[
 \left\lceil\frac{n-1}{3}\right\rceil
 =\left\lceil\frac n3\right\rceil,
\]
so there is nothing more to prove.

It remains to consider $n=3m+1$.  Suppose, contrary to the desired conclusion, that
\[
 d_P(p)\le m
 \qquad\text{for every }p\in P.
\]
The capacity bound gives the reverse inequality $d_P(p)\ge m$, hence
\[
 d_P(p)=m
 \qquad\text{for every }p\in P.
\]
For a fixed $p$, the $3m$ points of $P\setminus\{p\}$ are partitioned into exactly $m$ distance fibres, each of cardinality at most $3$.  Equality of the total cardinalities forces every one of those $m$ fibres to have cardinality exactly $3$.

Let
\[
 D:=\max\{|x-y|:x,y\in P\}
\]
and let $S\subseteq P$ be the set of all endpoints of pairs at distance $D$.  The set $S$ is nonempty and has at least two points.  If $p\in S$, then the distance $D$ is realized from $p$.  Since every distance fibre at $p$ has cardinality exactly $3$, the point $p$ has exactly three neighbours at distance $D$.  Every such neighbour also belongs to $S$.  Therefore the diameter graph on $S$ is $3$-regular.

The diameter of $S$ is still $D$: all distances in $S$ are at most $D$, and $S$ contains the two endpoints of at least one pair at distance $D$.  Hence \cref{thm:diameter}, applied to $S$, gives
\[
 |E(G_D(S))|\le|S|.
\]
On the other hand, the handshaking identity and $3$-regularity give
\[
 2|E(G_D(S))|=\sum_{p\in S}\deg(p)=3|S|,
\]
so
\[
 |E(G_D(S))|=\frac32|S|>|S|,
\]
a contradiction.  Thus some $p\in P$ satisfies $d_P(p)\ge m+1=\lceil n/3\rceil$.
\end{proof}

\begin{remark}
The diameter argument is needed only when $n\equiv1\pmod3$.  It rules out the hypothetical situation in which every distance circle about every point is filled to its maximum permitted capacity of three points.
\end{remark}

\section{An explicit two-axis construction}\label{sec:upper}

We now construct, for every $n$, an admissible configuration in which no point determines more than approximately $3n/4$ distances.

For an integer $r\ge0$, define
\[
 E_r:=
 \begin{cases}
  \varnothing,&r=0,\\[1mm]
  \{\pm2,\pm2^2,\ldots,\pm2^{r/2}\},&r\ge2\text{ even},\\[1mm]
  \{1\}\cup\{\pm2,\pm2^2,\ldots,\pm2^{(r-1)/2}\},&r\text{ odd}.
 \end{cases}
\]
Then $|E_r|=r$, every element is nonzero, and $E_r$ is a disjoint union of $\lfloor r/2\rfloor$ opposite pairs together with one unpaired element when $r$ is odd.  In particular, the number of orbits of $E_r$ under $t\mapsto-t$ is
\[
 \left\lceil\frac r2\right\rceil.
\]
Moreover, the product of any two elements of $E_r$ has the form $\pm2^k$ for some integer $k\ge0$.

Given $n\ge2$, put
\[
 a:=\left\lfloor\frac n2\right\rfloor,
 \qquad
 b:=\left\lceil\frac n2\right\rceil,
\]
and define
\[
 P_n:=\bigl(E_a\times\{0\}\bigr)
      \cup
      \bigl(\{0\}\times\sqrt3\,E_b\bigr).
\]
Because neither $E_a$ nor $E_b$ contains $0$, the two displayed parts are disjoint, and therefore $|P_n|=a+b=n$.

The next elementary criterion is the reason for the factor $\sqrt3$.

\begin{lemma}[Four points on two perpendicular axes]\label{lem:circlecriterion}
Let $u_1,u_2,v_1,v_2$ be nonzero real numbers with $u_1\ne u_2$ and $v_1\ne v_2$.  If the four points
\[
 (u_1,0),\quad (u_2,0),\quad (0,v_1),\quad (0,v_2)
\]
are concyclic, then
\[
 u_1u_2=v_1v_2.
\]
In fact, the equality is also sufficient.
\end{lemma}

\begin{proof}
Every Euclidean circle has an equation
\[
 x^2+y^2+\alpha x+\beta y+\gamma=0
\]
for suitable real coefficients $\alpha,\beta,\gamma$.  Restricting this equation to the $x$-axis shows that $u_1,u_2$ are the two roots of
\[
 t^2+\alpha t+\gamma=0.
\]
Vieta's formula gives $u_1u_2=\gamma$.  Restricting to the $y$-axis similarly gives $v_1v_2=\gamma$.  Hence $u_1u_2=v_1v_2$.

Conversely, if $u_1u_2=v_1v_2=\gamma$, then
\[
 x^2+y^2-(u_1+u_2)x-(v_1+v_2)y+\gamma=0
\]
passes through all four points.  Because it passes through two distinct real points, it represents a genuine circle rather than an empty or point circle.  Thus the equality is sufficient as well.
\end{proof}

\begin{proposition}[Admissibility of the construction]\label{prop:admissible}
No four points of $P_n$ are concyclic.
\end{proposition}

\begin{proof}
A circle intersects a fixed line in at most two points, because restricting its equation to that line gives a nonzero quadratic polynomial.  Hence a circle containing four points of $P_n$ would have to contain exactly two points on the $x$-axis and exactly two points on the $y$-axis.

Write their nonzero coordinates as $u_1,u_2\in E_a$ and
\[
 v_1=\sqrt3\,w_1,\qquad v_2=\sqrt3\,w_2,
 \qquad w_1,w_2\in E_b.
\]
By \cref{lem:circlecriterion}, concyclicity would imply
\[
 u_1u_2=v_1v_2=3w_1w_2.
\]
The left-hand side has the form $\pm2^r$, while the right-hand side has the form $\pm3\cdot2^s$, for some integers $r,s\ge0$.  Such numbers cannot be equal.  If their signs differ, equality is impossible.  If their signs agree, their absolute values are unequal because $2^r$ is not divisible by $3$, whereas $3\cdot2^s$ is divisible by $3$.  This contradiction proves the claim.
\end{proof}

\begin{proposition}[Distance count in the construction]\label{prop:count}
For $p\in E_a\times\{0\}$,
\[
 d_{P_n}(p)\le a-1+\left\lceil\frac b2\right\rceil.
\]
For $q\in\{0\}\times\sqrt3 E_b$,
\[
 d_{P_n}(q)\le b-1+\left\lceil\frac a2\right\rceil.
\]
\end{proposition}

\begin{proof}
Let $p=(u,0)$ with $u\in E_a$.  The other $a-1$ points on the $x$-axis contribute at most $a-1$ distinct distances.  For a point $(0,v)$ on the $y$-axis,
\[
 |(u,0)-(0,v)|=\sqrt{u^2+v^2}.
\]
This value is unchanged when $v$ is replaced by $-v$.  The set $\sqrt3 E_b$ has exactly $\lceil b/2\rceil$ orbits under $v\mapsto-v$, so all $b$ points on the $y$-axis contribute at most $\lceil b/2\rceil$ distinct distances from $p$.  Taking the union of the same-axis and cross-axis distance sets can only decrease the sum of these two upper bounds.  Therefore
\[
 d_{P_n}(p)\le a-1+\left\lceil\frac b2\right\rceil.
\]
The second inequality is identical after interchanging the two axes.
\end{proof}

\begin{theorem}[Upper bound]\label{thm:upper}
For every $n\ge2$,
\[
 f(n)\le
 \left\lfloor\frac{3n}{4}\right\rfloor-
 \ind_{\{4\mid n\}}.
\]
In particular,
\[
 f(n)\le\left\lfloor\frac{3n}{4}\right\rfloor,
\]
and for every $m\ge1$,
\[
 f(4m)\le3m-1.
\]
\end{theorem}

\begin{proof}
The configuration $P_n$ is admissible by \cref{prop:admissible}.  By \cref{prop:count}, its maximum number of distances is at most
\[
 \max\left\{
 a-1+\left\lceil\frac b2\right\rceil,
 b-1+\left\lceil\frac a2\right\rceil
 \right\},
 \qquad
 a=\left\lfloor\frac n2\right\rfloor,
 \quad
 b=\left\lceil\frac n2\right\rceil.
\]
The four residue classes of $n$ give the following values.
\begin{center}
\begin{tabular}{c@{\qquad}c@{\qquad}c@{\qquad}c@{\qquad}c}
\toprule
$n$ & $a$ & $b$
& $a-1+\lceil b/2\rceil$
& $b-1+\lceil a/2\rceil$\\
\midrule
$4m$   & $2m$   & $2m$   & $3m-1$ & $3m-1$\\
$4m+1$ & $2m$   & $2m+1$ & $3m$   & $3m$\\
$4m+2$ & $2m+1$ & $2m+1$ & $3m+1$ & $3m+1$\\
$4m+3$ & $2m+1$ & $2m+2$ & $3m+1$ & $3m+2$\\
\bottomrule
\end{tabular}
\end{center}
Thus the displayed maximum equals
\[
 \begin{cases}
 3m-1,&n=4m,\\
 3m,&n=4m+1,\\
 3m+1,&n=4m+2,\\
 3m+2,&n=4m+3,
 \end{cases}
\]
which is precisely $\lfloor3n/4\rfloor-\ind_{\{4\mid n\}}$.  Since $f(n)$ is the minimum over all admissible configurations, the same quantity is an upper bound for $f(n)$.
\end{proof}

Combining \cref{thm:lower,thm:upper} proves \cref{thm:main}.

\begin{remark}
The construction deliberately uses many collinear points.  This is permitted: the hypothesis excludes four concyclic points but imposes no restriction on collinear subsets.
\end{remark}

\section{The remaining linear-gap question}\label{sec:gap}

The capacity bound treats every occupied circle centred at $p$ as though it could always be filled with three points.  The following exact identity measures the failure of that idealized packing.

\begin{definition}
For $p\in P$ and $j\in\{1,2,3\}$, let $N_j(p)$ be the number of realized distances $r$ for which
\[
 |F_p(r)|=j.
\]
Define the distance-fibre deficiency at $p$ by
\[
 \Delta_P(p):=2N_1(p)+N_2(p).
\]
\end{definition}

\begin{proposition}[Exact deficiency identity]\label{prop:deficiency}
For every $p\in P$,
\[
 \boxed{
 3d_P(p)-(n-1)=\Delta_P(p).}
\]
Equivalently,
\[
 d_P(p)=\frac{n-1+\Delta_P(p)}{3}.
\]
\end{proposition}

\begin{proof}
The distance fibres partition $P\setminus\{p\}$, and every fibre has size $1$, $2$, or $3$.  Therefore
\[
 n-1=N_1(p)+2N_2(p)+3N_3(p),
\]
whereas
\[
 d_P(p)=N_1(p)+N_2(p)+N_3(p).
\]
Subtracting the first identity from three times the second gives
\[
 3d_P(p)-(n-1)=2N_1(p)+N_2(p)=\Delta_P(p).
\]
\end{proof}

\begin{corollary}[Exact reformulation of a linear improvement]\label{cor:reformulation}
Fix $c>0$.  For a given admissible $n$-point configuration $P$, the assertion that some $p\in P$ satisfies
\[
 d_P(p)>\left(\frac13+c\right)n
\]
is equivalent to
\[
 \max_{p\in P}\Delta_P(p)>3cn+1.
\]
Consequently, a universal estimate
\[
 f(n)>\left(\frac13+c\right)n
\]
for all sufficiently large $n$ is equivalent to proving that every sufficiently large admissible configuration has a point whose distance-fibre deficiency is larger than $3cn+1$.
\end{corollary}

\begin{proof}
By \cref{prop:deficiency},
\begin{align*}
 d_P(p)>\left(\frac13+c\right)n
 &\iff n-1+\Delta_P(p)>n+3cn\\
 &\iff \Delta_P(p)>3cn+1.
\end{align*}
Taking the maximum over $p\in P$ proves the first assertion.  Taking the minimum over all admissible configurations proves the second.
\end{proof}

Thus the second question in the problem is not decided by the present bounds.  A positive constant improvement over $1/3$ requires a genuinely global argument showing that, for at least one common centre, a positive proportion of the available three-point capacities remain unused.  The diameter-graph argument supplies only the constant-sized improvement needed in the residue class $n\equiv1\pmod3$.

\section{Small cases and final summary}

The first few exact values follow immediately:
\[
 f(1)=0,\qquad f(2)=1,\qquad f(3)=1,\qquad f(4)=2.
\]
Indeed, the first two are immediate; an equilateral triangle gives $f(3)=1$; and \cref{thm:main} gives $2\le f(4)\le2$.

The unconditional conclusions established in this note are therefore
\[
 \boxed{
 \left\lceil\frac n3\right\rceil
 \le f(n)
 \le
 \left\lfloor\frac{3n}{4}\right\rfloor-
 \ind_{\{4\mid n\}}.}
\]
The upper construction disproves the near-linear conjecture $f(n)>(1-o(1))n$.  Whether the lower bound admits a uniform improvement to $(1/3+c)n$ for some absolute $c>0$ is not resolved by these arguments; \cref{cor:reformulation} records the exact structural statement that such an improvement would require.

\end{document}
